\scanpage{161}
\mspage{106}
\unsure{n'est|n'ait} pas seulement une structure de \inserted{(soi-disant)} ``chemins'' \inserted{dans $X$} amorphe, \deleted{\unsure{avec}} seulement la donnée de la source et du but des chemins et des chemins triviaux (ces structures \inserted{sur $\mathrm{Ch}(X)$} dépendant fonctoriellement en $X$), mais qu'on a une composition des chemins, quand \deleted{\unsure{ceux}} ceux-ci sont ``composables'', i.e. une \deleted{\unsure{application}} flèche d'un produit fibré
\[
(c,c')\longmapsto c'\circ c:(\mathrm{Ch}(X),\beta_X)\times_X(\mathrm{Ch}(X),\sigma_X)\xrightarrow{\pi_X}\mathrm{Ch}(X)\text{,}
\]
i.e. on peut composer \unsure{ses|des} chemins sur \deleted{\unsure{X}} $X(T)$
\diagram{\begin{tikzcd}[ampersand replacement=\&]
x\arrow[r,"c"] \& y\arrow[r,"c'"] \& z
\end{tikzcd}}
ssi $\mathrm{but}\ c=\mathrm{source}\ c'$, et de plus on exige bien sûr que $\mathrm{source}\,(c'\circ c)=\mathrm{source}\ c$, $\mathrm{but}\,(c'\circ c)=\mathrm{but}\ c'$, ce qui s'exprime par un diagramme commutatif
\diagram{\begin{tikzcd}[ampersand replacement=\&]
(\mathrm{Ch}(X),\beta_X)\times_X(\mathrm{Ch}(X),\sigma_X)\arrow[rr]\arrow[dr] \& {} \& \mathrm{Ch}(X)\arrow[dl,"p_X"]\\
{} \& X\times X \& \text{.}
\end{tikzcd}}

La pensée naturelle, c'est qu'avec cette composition et ces ``\unsure{identités|unités}'', les $X(T)$, $\mathrm{Ch}(T)$ deviennent les objets et les flèches d'une \uline{catégorie} --- i.e. que $(X,\mathrm{Ch}(X),\deleted{\unsure{\pi_X}}\deleted{\unsure{\sigma}}\alpha_X,p_X,\pi_X)$ définissent une structure de ``catégorie dans $\mathcal M$''. Mais pour ce qu'on veut faire \inserted{maintenant} on n'a sans doute pas besoin de tout, \deleted{\unsure{de}} \unsure{cela|sla}.
% archive sheet number: 160
